\documentclass[11pt]{article}
\usepackage[margin=1in]{geometry}
\usepackage{xcolor}
\usepackage{booktabs}
\usepackage{tabularx}
\usepackage{siunitx}
\usepackage[hidelinks]{hyperref}
\definecolor{accent}{HTML}{334E68}
\setlength{\parindent}{0pt}
\setlength{\parskip}{0.5em}

\begin{document}

% ---------- Title page ----------
\begin{titlepage}
\begin{center}

{\color{accent}\rule{\textwidth}{2.5pt}}\\[0.4em]
{\large\scshape Meridian Structural Consultants}\\
{\small Civil and Structural Engineering \quad | \quad Portsmouth, UK}\\[0.2em]
{\color{accent}\rule{\textwidth}{0.8pt}}

\vspace{3.2cm}

{\Huge\bfseries\color{accent} Load Rating Analysis of the\\[0.25em] Alder Creek Road Bridge}\\[1.2em]
{\large Assessment of Live-Load Capacity under BD 21 Criteria}

\vspace{1.6cm}

\begin{tabular}{r l}
\textbf{Report No.} & MSC-TR-2026-041 \\
\textbf{Revision} & C \\
\textbf{Date} & 15 May 2026 \\
\textbf{Prepared for} & Wexford County Highways Department \\
\textbf{Prepared by} & R. Adeyemi, CEng \\
\textbf{Checked by} & L. Marchetti, CEng \\
\textbf{Status} & Issued for Review \\
\end{tabular}

\vfill

\textbf{\color{accent} Revision History}\\[0.4em]
\begin{tabularx}{\textwidth}{c c X c}
\toprule
\textbf{Rev} & \textbf{Date} & \textbf{Description} & \textbf{Approved} \\
\midrule
A & 20 Mar 2026 & Initial draft for internal review & LM \\
B & 14 Apr 2026 & Updated girder section losses from site inspection & LM \\
C & 15 May 2026 & Incorporated client comments, revised load factors & LM \\
\bottomrule
\end{tabularx}

\vspace{1em}
{\color{accent}\rule{\textwidth}{0.8pt}}\\[0.3em]
{\small This document is confidential and prepared solely for the named client.
It may not be reproduced without written permission.}

\end{center}
\end{titlepage}

% ---------- Body ----------
\section{Introduction}
The Alder Creek Road Bridge is a two-span steel girder bridge built in 1968
carrying a single carriageway over Alder Creek. Wexford County Highways
Department commissioned this assessment after routine inspection identified
section loss in the bottom flanges of both main girders.

\section{Scope}
This report determines the safe live-load capacity of the superstructure,
accounting for measured corrosion losses, and recommends any interim traffic
restrictions. Substructure and hydraulic assessments are outside the scope.

\section{Summary of Findings}
The governing girder retains \SI{87}{\percent} of its original flexural
capacity. The bridge can safely carry vehicles up to \SI{32}{\tonne} gross
weight without restriction. A repainting and monitoring programme is
recommended within 24 months to arrest further section loss.

\end{document}
